\section{Выводы}

Робот написан и работает: он собрал корпус из 79\,566 документов с двух
источников --- русской Википедии и film.ru, --- пережил остановки, сон и
перезагрузку ноутбука и обрывы сети, а переобкачка находит и обновляет
изменившиеся документы. Переобкачка вынесена в отдельную программу и связана с
роботом через очередь.

Самое полезное, что дала работа, --- это несколько ошибок, которые на бумаге не
видны.

\textbf{Фильтр \enquote{что обходить} и фильтр \enquote{что сохранять} --- это
разные фильтры.} В первой версии робот складывал в базу оглавления и каталоги:
текста в них достаточно, чтобы пройти проверку по длине, и формально они
проходили все условия. Понадобился отдельный набор правил, решающий, что
считать документом, а что --- только источником ссылок.

\textbf{Без ограничения на то, откуда брать ссылки, корпус перестаёт быть
тематическим.} Из статьи о фильме ссылки ведут на статьи о городах, годах и
профессиях, и через несколько шагов обхода робот собирает уже не кино. Пришлось
разрешить сбор ссылок только со страниц категорий и каталогов.

\textbf{Стандартная библиотека умеет подвести молча.} Робот отказывался качать
Википедию целиком, помечая все адреса как запрещённые. Причина оказалась в том,
что \texttt{RobotFileParser.read()} ходит за \texttt{robots.txt} через
\texttt{urllib} со стандартным User-Agent, Викимедиа отвечает на него
\texttt{403}, а парсер трактует \texttt{403} как \enquote{запрещено всё}.
Исправление --- качать \texttt{robots.txt} той же сессией, что и остальные
страницы, и разбирать его через \texttt{parse()}.

\textbf{Адрес нужно нормализовать раньше, чем кажется.} Ссылки в HTML Википедии
процентно закодированы, и пока адреса не раскрывались, правила с кириллицей по
ним не срабатывали: в базу лезли служебные страницы, которые фильтр должен был
отсечь.

\textbf{Сотня документов не показывает того, что покажут десятки тысяч.}
Демонстрационный прогон прошёл гладко, а полная обкачка вскрыла сразу
несколько проблем: продолжение категорий закрыто \texttt{robots.txt}, через
родительские категории робот уходит из тематики, тысячи красных ссылок тратят
запросы впустую, а защита второго источника включается только после сотен
запросов. Проверять источник и правила обхода надо на масштабе, а не на образцах.

\textbf{Сеть ненадёжна, и робот должен это учитывать.} Ошибка соединения не
означает, что страницы нет. Пока робот считал любую ошибку окончательной, один
ночной обрыв интернета стоил двенадцати тысяч адресов. Временные сбои нужно
отличать от настоящих ошибок и повторять с паузой.

\textbf{Второй источник надо выбирать с запасом.} Кино-Театр.Ру по пяти
страницам выглядел подходящим, но закрылся капчей после пятисот запросов, и
блокировка не снималась сутками. Замена на film.ru оказалась быстрой
только потому, что робот с самого начала описывал источник конфигом, а не
кодом.

\textbf{Проверять выделение текста нужно глазами и тестами, а не по
счётчикам.} Робот отчитывался о сохранённых документах, объёмы выглядели
правдоподобно, и то, что у второго источника в текст попадает вся страница,
обнаружилось только тестом на настоящей странице: в тексте статьи нашлась
строка из подвала сайта.

\subsection*{Недостатки}

Робот однопоточный. При паузах в одну и две секунды предел --- около полутора
запросов в секунду на два источника, а на деле выходит около 0,5: остальное
съедает ожидание ответа. Полная обкачка заняла четыре дня. Выбор источника по
готовности хоста убирает простой, но не больше: настоящее решение --- несколько
потоков с очередью на каждый хост, как в схеме URL frontier с лекции.

Корпус Википедии собран неравномерно: из каждой категории берутся только первые
200 статей по алфавиту, так что статьи в начале алфавита представлены лучше.
Полный список статей категории без \texttt{/w/} не получить; честный путь ---
дамп Википедии или её API, но это уже не обход страниц.

Около 0,8\,\% документов Википедии не относятся к кино: категории о фильмах
иногда включают основную статью темы (\enquote{Канада}, \enquote{Франция}).
Отфильтровать их можно по списку категорий самой статьи, который есть в её
HTML, --- так эта доля и посчитана.

Около пятой части запросов уходит на страницы категорий и списков, которые
документами не становятся.

\enquote{Сырой} HTML занимает на порядок больше места, чем текст. MongoDB
по умолчанию сжимает данные, и на диске база втрое меньше их объёма, но
явное сжатие текста словарными методами дало бы больше.

Переобкачка обходит документы по времени последней проверки, то есть одинаково
часто для всех. Разумнее было бы учитывать, как часто документ менялся раньше, и
проверять стабильные страницы реже, а живые --- чаще.

Наконец, робот не защищён от ловушек: календари и страницы с бесконечной
пагинацией порождают бесконечную очередь. Пока от этого спасают только
ограничение глубины и фильтры адресов.

\pagebreak
